% Part: sets-functions-relations
% Chapter: induction
% Section: induction

\documentclass[../../include/open-logic-section]{subfiles}

\begin{document}

\olfileid[hi]{sfr}{ind}{int}

\olsection{परिचय}

\begin{explain}
आगमन, तर्कशास्त्र में गणित की सामान्यतः प्रयुक्त प्रमाण-तकनीक है। आप इसे
गणित से जानते होंगे, जहाँ प्राकृतिक संख्याओं पर आगमन किया जाता है। वास्तव
में अनेक भिन्न प्रकार के गणितीय और तार्किक वस्तु-समुच्चय आगमन के अनुकूल होते
हैं।

सामान्यतः आगमन हमें सार्वत्रिक दावा सिद्ध करने देता है; अर्थात् यह दिखाना कि
किसी निश्चित प्रकार की प्रत्येक वस्तु में कोई गुण है। विशेषतः आगमन प्रायः तब
उपयोगी है जब प्रमाण की ``सार्वत्रिक प्रस्तावना'' विधि बहुत जटिल हो। आगमन केवल
विशेष ढंग से निर्मित गणितीय वस्तुओं पर काम करता है: यदि समुच्चय का प्रत्येक
अवयव या तो मूल है अथवा मूल अवयवों से बना है, तो हम उस पर आगमन का उपयोग कर
सकते हैं। अधिक ठीक-ठीक:
\end{explain}

\begin{defn}
कोई समुच्चय $S$, फलन $f$ के अंतर्गत \emph{संवृत} है तभी और केवल तभी जब
$x \in S$ होने पर $f(x) \in S$।
\end{defn}

\begin{explain}
उदाहरणतः प्राकृतिक संख्याएँ ($\mathbb{N}$), उत्तरवर्ती फलन $f(x) = x+1$ के
अंतर्गत संवृत हैं। यदि $x$ प्राकृतिक संख्या है, तो $x+1$ भी है। किंतु
प्राकृतिक संख्याएँ पूर्ववर्ती फलन $f(x) = x-1$ के अंतर्गत संवृत नहीं हैं,
क्योंकि $0$ प्राकृतिक संख्या है पर $-1$ नहीं।
\end{explain}

\begin{defn}
गुण $P$ को फलन $f(x)$ के \emph{अंतर्गत संरक्षित} कहते हैं, जब $f$ के प्रांत
की किसी भी वस्तु $a$ के लिए $P(a)$ से $P(f(a))$ निष्पन्न हो (यदि $a$ में
गुण $P$ है, तो $f(a)$ में भी है)।
\end{defn}

\begin{explain}
``सम संख्या होने'' का गुण फलन $f(x) = x+2$ के अंतर्गत संरक्षित है, क्योंकि
यदि $x$ सम है, तो $x+2$ भी सम है। सम संख्या होने का गुण उत्तरवर्ती फलन के
अंतर्गत संरक्षित नहीं है, क्योंकि यदि $x$ सम है, तो $x+1$ विषम है।

आगमन प्राकृतिक संख्याओं पर इसलिए काम करता है क्योंकि उत्तरवर्ती फलन को परिमित
बार लागू करके प्रत्येक प्राकृतिक संख्या तक $0$ से पहुँचा जा सकता है। यह
किसी भी ऐसे समुच्चय का अभिलक्षण है जिस पर हम आगमन कर सकते हैं।
\end{explain}

\begin{defn}[$\Nat$ पर आगमन]
यदि 0 में कोई गुण $P$ है और $P$ उत्तरवर्ती फलन के अंतर्गत संरक्षित है, तो
प्रत्येक प्राकृतिक संख्या में गुण $P$ है।
\end{defn}

\begin{ex} पहली $n$ प्राकृतिक संख्याओं का योग $\frac{n(n+1)}{2}$ है। \\

\textbf{आधार स्थिति।} पहली 0 प्राकृतिक संख्याओं का योग $0=
\frac{0\cdot 1}{2}$ है।\\

\textbf{आगमन चरण।} मान लें कि गुण $k$ के लिए सत्य है। हम दिखाएँगे कि वह
$k+1$ के लिए सत्य है। दूसरे शब्दों में मान लें कि $P(k)$ सत्य है, अर्थात्
\[ 0 + 1 + \ldots + k = \frac{k(k+1)}{2} \]
हम दिखाएँगे कि $P(k+1)$ सत्य है, अर्थात्
\[ 0 + 1 + \ldots + k + (k+1) = \frac{(k+1)(k+2)}{2} \]
$P(k)$ की मान्यता का उपयोग करके।

\begin{align*}
0 + 1 + \ldots + k &= \frac{k(k+1)}{2} \tag{मान्यता} \\
0 + 1 + \ldots + k + (k+1) &= \frac{k(k+1)}{2} + (k+1)
\tag{दोनों पक्षों में $k+1$ जोड़ें} \\
&= \frac{k(k+1) + 2(k+1)}{2} \\
&= \frac{k^2 + k + 2k +2}{2} \\
&=\frac{(k+1)(k+2)}{2}
\end{align*}
अतः आगमन चरण पूरा होता है और हमने दावा सिद्ध कर दिया।\footnote{अनुवादक की
टिप्पणी: जमी हुई स्रोत प्रति में अंश का पहला पद $2k$ छपा है। ठीक पिछली
पंक्ति में $k(k+1)$ का विस्तार $k^2+k$ है; यहाँ उसी के अनुसार $k^2$ किया गया है।}
\end{ex}

\begin{explain}
!!{formula}s के लिए मूल अवयव परमाण्विक !!{formula}s हैं। संकारकों से कम
जटिल !!{formula}s को जोड़कर अधिक जटिल !!{formula}s मिलते हैं। प्रत्येक
संकारक को !!{formula}s के ऐसे फलन के अनुरूप माना जा सकता है जो संबद्ध संकारक
से दोनों को जोड़ने वाला नया !!{formula} देता है। उदाहरणतः दो !!{formula}s
$!A$ और $!B$ को संयोजन में जोड़ने वाले फलन को $\mathcal E _\land (!A,
!B) = !A \land !B$ लिख सकते हैं। इन्हें हम \emph{सूत्र-निर्माण फलन} कह
सकते हैं। !!{formula}s का समुच्चय इन फलनों के अंतर्गत संवृत है: जब भी $!A$
और $!B$ !!{formula}s हों, तब $\lnot !A$, $!A \land !B$ आदि भी हैं।\footnote{अनुवादकीय
टिप्पणी: स्रोत के इस परिच्छेद और नीचे की आगमन-विधि में सूत्र के लिए दिया
विकृत नियंत्रण-संकेत संकलित नहीं होता। यहाँ उसे सूत्र-संकेत $!A$ किया गया है;
मूल स्रोत-पाठ और परिवर्तन का विवरण साथ की उत्पत्ति-अभिलेख सामग्री में सुरक्षित हैं।}
\end{explain}

\begin{defn}[!!{formula}s पर आगमन]
यदि प्रत्येक परमाण्विक !!{formula} में गुण $P$ है और $P$, !!{formula}-निर्माण
फलनों के अंतर्गत संरक्षित है, तो प्रत्येक !!{formula} में गुण $P$ है।
\end{defn}

\begin{explain}
यह परिभाषा !!{formula}s पर आगमनात्मक प्रमाणों की एक ``विधि'' सुझाती है। यदि
हमें दिखाने को कहा जाए कि प्रत्येक !!{formula} में गुण $P$ है, तो निम्न चरणों
का अनुसरण करें:\\

\textbf{आधार स्थिति।} मान लें $!A$ कोई परमाण्विक !!{formula} है। [\ldots]
अतः $!A$ में गुण $P$ है।

\textbf{आगमन चरण।} मान लें $!A$ और $!B$ ऐसे !!{formula}s हैं जिन दोनों
में गुण $P$ है।

स्थिति $\lnot$: [\ldots] अतः $\lnot !A$ में गुण $P$ है।

स्थिति $\land$: [\ldots] अतः $!A \land !B$ में गुण $P$ है।

स्थिति $\lor$: [\ldots] अतः $!A \lor !B$ में गुण $P$ है।

स्थिति $\lif$: [\ldots] अतः $!A \lif !B$ में गुण $P$ है।

स्थिति $\liff$: [\ldots] अतः $!A \liff !B$ में गुण $P$ है।

स्थिति $\lforall[]$: [\ldots] अतः $\lforall[x] !A$ में गुण $P$ है।

स्थिति $\lexists[]$: [\ldots] अतः $\lexists[x] !A$ में गुण $P$ है। \\

अतः प्रत्येक !!{formula} में गुण $P$ है।
\end{explain}

\end{document}
